\documentclass[10pt,a4paper]{amsart}
\usepackage[margin=19mm]{geometry}
\usepackage{amsmath,amssymb,mathtools}
\usepackage[scr=rsfs,cal=euler]{mathalfa}
\usepackage{xfrac}
\usepackage[unicode,hidelinks,bookmarks=false]{hyperref}
\newcommand{\ie}{i.e.\ }
\newcommand{\cf}{cf.\ }
\newcommand{\resp}{respectively\ }
\newcommand{\Subsection}{\subsection}
\providecommand{\nobreakdash}{\nobreak\hbox{-}}
\setlength{\emergencystretch}{2em}
\multlinegap0pt
\usepackage{amssymb}
\usepackage{mathtools}
\usepackage{tikz}
\usepackage{xcolor}
\newtheorem{lemma}{Lemma}
\title{Long-range expanders: construction and cutoff}
\author{Dylan J. Altschuler}
\date{Source: arXiv:2609.14945v1 (CC BY 4.0)}
\begin{document}
\maketitle

\noindent\emph{Source and licence.} Long-range expanders: construction and cutoff. Version arXiv:2609.14945v1 (CC BY 4.0), distributed by arXiv under the Creative Commons Attribution 4.0 International licence, http://creativecommons.org/licenses/by/4.0/, as recorded in the arXiv licence metadata for this exact version and shown on the abstract page https://arxiv.org/abs/2609.14945v1. Full licence notice: \texttt{licenses/arxiv-2609.14945-CC-BY-4.0.txt}.\par\smallskip

\noindent\textit{Editorial scope.} Counted unit: the article's lemma lem:directedentropy (source0.tex lines 590--603). The article's complete original proof of it is delivered with it (source0.tex lines 605--717, one \texttt{proof} environment of 113 lines). No mathematics is written, repaired or re-derived: the statement and the proof are the article's own lines, in the article's own notation, and no retained line is substituted or re-flowed.  No further source-local context is needed: every cross-reference of every retained line resolves inside the two delivered spans.  Nothing was omitted from any retained span, so the package carries no omission record.  No retained line contains a citation, so the wrapper carries no bibliography and no bibliography entry.  No retained line contains a figure environment, an image-inclusion command or drawing code, so no image is delivered and the package declares no asset dependency.  Editorial formatting only, in addition: an AMS \texttt{amsart} preamble that restates the article's own declarations and macros used by the retained lines, namely the environments \texttt{lemma} on separate counters, together with the standard packages those lines need.  The retained environments print in this document as Lemma 1 in order of appearance, whereas the article prints the same items with the numbering of its own full document.  The complete native source of the article as distributed in the arXiv e-print for this exact version is bundled unchanged as \texttt{sources/210375/source0.tex}.
\begin{lemma}\label{lem:directedentropy}
Let $\Omega$ be the vertex set of a directed graph on $N$ states,
with $q\ge2$ incoming and $q$ outgoing edges at every state. Suppose that for all $\emptyset \neq S \subseteq \Omega$ and every integer $r\ge0$,
\[
 |B(S,r)|\ge a\min\{N,|S|q^r\}.
\]
Then, there are constants $C_0\ge1$ and $\delta>0$, depending only on
$a,q$, such that a simple random walk $(Y_t)_{t\ge0}$ started uniformly on
$\emptyset \neq S \subseteq \Omega$ satisfies, for every integer $h\ge0$,
\begin{equation}\label{eq:entropybound}
 H(Y_h)\ge\log|S|+h\log q-\delta^{-1}\log(qh+1)
 \qquad\text{if }h+C_0\le\log_q(N/|S|).
\end{equation}
\end{lemma}

\begin{proof}
Assume $a\le1$, choose an integer $C_0\ge1$ with
$q^{-C_0}\le a(q-1)/2$, and put $\delta=a(q-1)/(2q)$.
We induct on $h$, simultaneously for all eligible starting sets $S$.
The assertion is immediate when $h=0$.

Write $\Gamma_r(S)$ for the states at directed distance exactly $r$
from $S$. Each state of $\Gamma_{r+1}(S)$ has a predecessor in
$\Gamma_r(S)$, so $|\Gamma_{r+1}(S)|\le q|\Gamma_r(S)|$.
Using growth at radius $h+C_0$, we obtain
\[
 \begin{aligned}
 a|S|q^{h+C_0} \le |B(S,h+C_0)| 
 &\le |S|\sum_{r=0}^{h-1}q^r
       +|\Gamma_h(S)|\sum_{r=0}^{C_0}q^r \le\frac{|S|q^h+|\Gamma_h(S)|q^{C_0+1}}{q-1}.
 \end{aligned}
\]
Thus
\begin{equation}\label{eq:layer}
 |\Gamma_h(S)|\ge\delta|S|q^h.
\end{equation}

For each state in $\Gamma_r(S)$, $r\ge1$, choose one incoming edge
from $\Gamma_{r-1}(S)$. The chosen edges form a forest rooted at $S$:
there is exactly one forest path from $S$ to each reachable state.
Let $T$ be the first time the walk uses an edge outside this forest.
All length-$h$ trajectories have probability $1/(|S|q^h)$, so
\begin{equation}\label{eq:success}
 \alpha :=\Pr(T>h)=\frac{|\Gamma_h(S)|}{|S|q^h}\ge\delta,
 \qquad \text{ and } \qquad
 \mathcal L(Y_h\mid T>h)=\operatorname{Unif}(\Gamma_h(S)).
\end{equation}

On departure from the forest, conditioning only on $\{T=t\}$ need not
leave a uniform law: several departing edges may have the same
endpoint. To distinguish them, list the nonforest edges from
$\Gamma_{t-1}(S)$ into each state $y$ in an arbitrary fixed order.
Let $m_t(y)\le q$ be their number, and let
\[
 S_{t,i}=\{y:m_t(y)\ge i\},\qquad 1\le i\le q.
\]
We consider $1\le t\le h$ and omit pairs $(t,i)$ of probability zero.
On $\{T=t\}$, let $I$ be the index of the edge used in the list at
$Y_t$. For each $y\in S_{t,i}$, exactly one trajectory realizes
$\{T=t,I=i,Y_t=y\}$: follow the forest to the selected edge and
then traverse it. Therefore
\begin{equation}\label{eq:branch}
 \beta_{t,i}:=\Pr(T=t,I=i)=\frac{|S_{t,i}|}{|S|q^t},
 \qquad \text{ and } \qquad
 \mathcal L(Y_t\mid T=t,I=i)=\operatorname{Unif}(S_{t,i}).
\end{equation}
This conditioning involves only the first
$t$ steps, so the remaining walk has its usual transition law.
Moreover, $|S_{t,i}|\le|S|q^t$ implies
\[
 h-t+C_0\le\log_q(N/|S|)-t\le\log_q(N/|S_{t,i}|).
\]
The induction hypothesis consequently applies to the remaining
$h-t$ steps, giving
\[
 H(Y_h\mid T=t,I=i)
 \ge \log|S|+h\log q+\log\beta_{t,i}
      -\delta^{-1}\log(q(h-t)+1).
\]
On the event $\{T>h\}$, \eqref{eq:success} gives the conditional entropy of $Y_h$ as
$\log|S|+h\log q+\log\alpha$.

Let $Z$ indicate whether the walk stays in the forest through
time $h$, and otherwise specify its departure time and edge index:
\[
 Z=
 \begin{cases}
  0,&T>h,\\
  (T,I),&T\le h.
 \end{cases}
\]
Thus $\Pr(Z=0)=\alpha$ and $\Pr(Z=(t,i))=\beta_{t,i}$.
Since $Z$ takes at most $qh+1$ values,
\[
 H(Z)=-\alpha\log\alpha-\sum_{t,i}\beta_{t,i}\log\beta_{t,i}
 \le\log(qh+1).
\]

Conditioning on $Z$ separates the walk into the cases
estimated above. By concavity of entropy,
\[
 H(Y_h)
 \ge H(Y_h\mid Z) =\alpha H(Y_h\mid T>h)
   +\sum_{t,i}\beta_{t,i}H(Y_h\mid T=t,I=i).
\]
Substituting the preceding conditional entropy bounds, and using
$\alpha+\sum_{t,i}\beta_{t,i}=1$, gives

\begin{align*}
 H(Y_h)
 &\ge \log|S|+h\log q
      +\alpha\log\alpha+\sum_{t,i}\beta_{t,i}\log\beta_{t,i}
      -\delta^{-1}\sum_{t,i}\beta_{t,i}\log(q(h-t)+1)\\
 &=\log|S|+h\log q-H(Z)
      -\delta^{-1}\sum_{t,i}\beta_{t,i}\log(q(h-t)+1)\,.
\end{align*}
Recall that $H(Z)\le\log(qh+1)$ and
$\sum_{t,i}\beta_{t,i}=1-\alpha\le1-\delta$. Using these relations and the crude inequality $\log(q(h-t)+1)\le\log(qh+1)$,
\[
 \begin{aligned}
 H(Y_h)  &\ge\log|S|+h\log q-\log(qh+1)
      -\frac{1-\delta}{\delta}\log(qh+1)\\
 &=\log|S|+h\log q-\delta^{-1}\log(qh+1).
 \end{aligned}
\]

This completes the induction.
\end{proof}

\end{document}
